\section{Introduction}\label{sec:intro}

A printed page makes two commitments at once. It commits to a linguistic content, the words and sentences that a reader is meant to recover, and it commits to a convention for turning that content into marks, the alphabet, the orthography and the particular letter shapes of a typeface. Ordinary reading hides the second commitment because the convention is learned so early and rehearsed so heavily that it ceases to be visible as a convention at all. A reader who meets a page set in an unfamiliar script is therefore in an unusual position: the content is still there, but the instrument that delivers it has to be learned first, and the usual pedagogical response is to teach the instrument separately, through an alphabet chart, a primer, or a course, before the reading proper begins.

The constructions studied in this essay start from a different question. Suppose the instrument could be changed while the content continued to be delivered, so that a reader never stopped to study a table of signs but met the new forms inside running text, a few at a time, with enough surrounding material in the familiar form to carry the argument. A book arranged in this way would drift. Its first pages would use forms the reader already knows, its later pages would introduce small departures, and its last pages would be set in a script that would have been hard to read on the first page. The book would supply part of the preparation needed to read its own ending. Whether it does so, and under what conditions, is an empirical matter, and the essay treats it as one. What can be done in advance, and what occupies most of the following pages, is to specify the objects involved precisely enough that the empirical question can be asked without ambiguity: what a visual substitution is, how resemblance between letters can be measured and optimised, how a display can be separated from a stored text, how the intermediate forms of a transition behave, how a schedule of changes interacts with the frequency of letters in running text, and what mathematical structures describe the resulting editions.

The starting point is a visual substitution cipher, a one-to-one replacement of letters in which each letter is exchanged for one that looks like it. The purpose is the reverse of the usual purpose of a cipher. The substitution is intended to be easy to see through, so that a reader gradually learns to read it without consciously decoding each letter, with spaces, punctuation and capitalisation preserved so that word shape and sentence context can support recognition where an individual replacement is ambiguous. Such an alphabet can be rendered in two quite different ways. A font can display each stored character exactly as the ordinary letter it represents, so that the stored text is substituted while the page looks normal, or it can display a shape somewhere between the stored character and its partner, so that the substitution is perceptible but the familiar structure remains available to the reader. A font of the first kind already exists in the project's font collection: the LingoJam Cipher font, which displays characters of a particular substitution cipher as ordinary Standard Galactic letters, so that the stored characters carry one representation and the font supplies the visual interpretation. The present essay generalises that separation and asks what happens when the font is not fixed but changes from one part of a book to another.

Three further motivations shape the treatment. The first is theoretical. The arbitrariness of the sign, in the sense that no necessity ties a linguistic unit to its particular written shape, is a familiar thesis, and a collection of working ciphers, typefaces and drifting editions is a way of exhibiting it in a form that can be read rather than merely asserted. Yet arbitrariness of assignment does not imply equal ease of acquisition, and the visual substitution alphabet makes the gap between the two concrete: the exchange of B and P can be learned through a perceptual bridge that the exchange of C and K lacks. The second motivation is practical. The essays in this project are released with their source in the public domain, so that readers can propose corrections and can also print any text in whatever script they choose. This turns each text into a family of editions, and the editions are at once publications and research instruments. The third is methodological. Because the editions are generated from separately identifiable components, the conditions of a transition can be varied one at a time, and the observations of how readers cope can be attached to a precise description of what they read.

\subsection{Status of the claims}

The essay combines results of different kinds, and the difference matters because the temptation in a subject of this sort is to let a plausible mechanism stand in for a measured outcome. Each major claim is therefore tagged below by what supports it. Statements of the first kind are proved from definitions given in the text. Statements of the second kind are computed, which means that they report the output of a program on stated inputs, the program being reproduced in the appendices, and they are exactly as strong as the inputs and the model behind them. Statements of the third kind are hypotheses about readers, and the essay does not claim to have established any of them.

\begin{table}[ht]
\centering
\small
\renewcommand{\arraystretch}{1.25}
\begin{tabular}{>{\raggedright\arraybackslash}p{2.7cm}>{\raggedright\arraybackslash}p{6.3cm}>{\raggedright\arraybackslash}p{2.7cm}}
\toprule
Section & Claim & Status\\
\midrule
\ref{sec:cipher} & Size of the space of reciprocal alphabets; round-trip recovery & Proved\\
\ref{sec:resemblance} & Optimal pairings differ across typefaces; overlap measure disagrees with intuition for several pairs & Computed\\
\ref{sec:drift} & Linear exchange of two letters passes through a glyph collision; a detour removes it; three-cycles avoid it in general position & Proved\\
\ref{sec:drift} & Exposure conservation and a lower bound on the length of a drift; rare-first ordering attains it in the cases tried & Proved (bound), computed (attainment)\\
\ref{sec:arabic} & Context-dependent display breaks the commuting of substitution and rendering unless joining classes are preserved & Proved\\
\ref{sec:formal} & Gluing criterion for local keys; the true local keys of a swap-ending drift are compatible but have no global key for any cover & Proved\\
\ref{sec:empirical} & Drifting presentation supports receptive acquisition and transfer & Hypothesis\\
\bottomrule
\end{tabular}
\caption{Claims of the essay by kind of support.}\label{tab:status}
\end{table}

\subsection{Plan}

Section~\ref{sec:cipher} defines visual substitution alphabets and counts them. Section~\ref{sec:resemblance} treats resemblance as a weight to be optimised and reports a computed comparison across typefaces. Section~\ref{sec:arb} sets the arbitrariness of the sign against the constraints on acquisition. Section~\ref{sec:edition} describes an edition as a composite of independently variable components and discusses the font technology that realises it. Section~\ref{sec:drift} analyses intermediate forms and schedules of change. Section~\ref{sec:arabic} considers contact between Latin, Arabic and Standard Galactic forms, with attention to terminology. Section~\ref{sec:source} discusses the publication of editable source and the provenance of editions. Section~\ref{sec:empirical} specifies an empirical programme. Section~\ref{sec:formal} gives formal interpretations in category theory, type theory and sheaf theory. Section~\ref{sec:concl} concludes, and the appendices contain the derivations, the full computational results and the code.
